\documentclass{report}
\usepackage{amsmath,amssymb}
\setlength{\parindent}{0mm}
\setlength{\parskip}{1em}
\begin{document}
\begin{center}
\rule{6in}{1pt} \
{\large Stefan Henneking \\
{\bf Evanescent data segmental refinement multigrid: ideas and experiments}}

100 10th Street NW \\ Apt 706 \\ Atlanta \\ GA 30309
\\
{\tt henneking@gatech.edu}\\
Richard Vuduc\\
Mark F. Adams\end{center}

We introduce a modified approach to evanescent data segmental refinement
(SR), based on Achi Brandt's work on low-memory multi-level techniques
(1977). Evanescent data SR avoids storing the entire finest grid(s) at
any time during the computation and thereby drastically reduces the
memory complexity of a standard full multigrid (FMG) full approximation
scheme solver. Naturally, this feature also provides a way to compress an
accurate fine-grid solution to a coarser grid and reproduce the fine-grid
solution later. We discuss a simple one-level evanescent data scheme as
well as approaches to a multi-level version; both can be applied within a
sequential SR solver or as part of a parallel one. We present preliminary
experimental results of a model problem for two FMG SR solvers: a
vertex-centered finite difference solver and a cell-centered finite
volume solver. We discuss the trade-off between storage requirements and
computational complexity.


\end{document}
